\scanpage{145}
\mspage{94}

Supposons en effet cette dernière condition satisfaite, et prouvons que $H$ est tot. $\mathcal W$-asphérique, i.e. que pour tout $p:U\unsure{-|\to}U'$, $U'/p_!(H)$ soit $\mathcal W$-asph.

\diagram{\begin{tikzcd}[ampersand replacement=\&]
V\overset{\text{déf}}{=}U_{/H}\arrow[r,"{\unsure{\alpha|p}}"]\arrow[dr,"f"'] \& U\arrow[d,"p"] \\
\& U'
\end{tikzcd}}

Soit $f$ le foncteur composé du diagramme ci-contre, on a $H\simeq\alpha_!(e_V)$ \deleted{\unsure{(par le|(par la}} (par $(*)$ p. 92), d'où
\[
p_!(H)\simeq p_!\alpha_!(e_V)\simeq(p\alpha)_!(e_V)\simeq f_!(e_V),
\]
qui est $\mathcal W$-asph. par hypothèse sur $V=U_{/H}$.

La réciproque de prop. 8 est-elle valable, i.e. si $H$ est totalement $\mathcal W$-asph., en est-il de \unsure{=|même} de $V=U_{/H}$ ? La difficulté provient du fait que \unsure{cela|cette} \deleted{\unsure{condition signifie|condition, signifie}} \inserted{signifie} que pour \uline{toute} flèche $f\inserted{{:}V\to V'}$ dans $\mathrm{Cat}$ de source $V$, on ait $f_!(e_V)$ $\mathcal W$-asphérique, alors que l'hypothèse \deleted{\unsure{ne pas|ne peut}} équivaut à le postuler seulement pour les flèches qui se factorisent par $U$. Notons que si on ne fait pas cette hypothèse sur $f$, mais \inserted{qu'on suppose} en lieu de cela que $f$ s'insère dans un carré commutatif

% archive sheet number: 145
